\subsubsection{次调和函数}\label{subsec:subharmonic-functions}

需要一个 Dirichlet 问题下解的概念, 不要求下解属于 $C^2(\Omega)$ 类. 这一定义由最大值原理启发而来.

\begin{Definition}
\label{def:continuous-subharmonic-function}
若对每个满足 $\overline B\subset\Omega$ 的球 $B$, 以及每个函数 $h\in C(\overline B)$, 只要 $h$ 在 $B$ 内调和且在 $\partial B$ 上满足 $u\leq h$ ($u\geq h$), 就有在 $B$ 内 $u\leq h$ ($u\geq h$), 则称 $C^0(\overline\Omega)$ 中的函数 $u$ 为\textbf{次调和的} (\textbf{超调和的}). Dirichlet 问题的下解 (上解) 是次调和 (超调和) 函数 $u\in C(\overline\Omega)$, 且在 $\partial\Omega$ 上 $u\leq g$ ($u\geq g$).
\end{Definition}

显然, 若 $\Delta u\geq0$, 则 $u$ 在新定义的意义下也是次调和的. 指出以下性质:

\begin{enumerate}
\item 强最大值原理成立, 即若 $u$ 次调和, $v$ 超调和, 且在 $\partial\Omega$ 上 $v\geq u$, 则或者在 $\Omega$ 内 $v>u$, 或者处处有 $v=u$. 用反证法证明. 假设 $u-v$ 在某点 $\bm{x}_0\in\Omega$ 取得最大值 $M$, 其中 $M\geq0$. 若整个 $\Omega$ 内 $u-v\equiv M$, 则有 $u=v$; 因而可假设 $\Omega$ 内存在点使 $u-v\not\equiv M$. 在这种情形下, 可以适当选择 $\bm{x}_0$, 使存在以 $\bm{x}_0$ 为球心的球 $B\subset\Omega$, 且 $u-v$ 在 $\partial B$ 上不处处等于 $M$. 设 $\bar u$ 和 $\bar v$ 为 $B$ 上的调和函数, 分别在 $\partial B$ 上等于 $u$ 和 $v$. 得到
\begin{equation}\label{eq:subharmonic-difference-at-maximum}
M=(u-v)(\bm{x}_0)\leq(\bar u-\bar v)(\bm{x}_0),
\end{equation}
而由调和函数的强最大值原理, 右端严格小于 $M$. 因而矛盾. 一个直接推论是, Dirichlet 问题的每个下解都小于或等于每个上解.\item 设 $u$ 在 $\Omega$ 内次调和, $B$ 为满足 $\overline B\subset\Omega$ 的球. 设 $\bar u$ 为 $B$ 上满足在 $\partial B$ 上 $\bar u=u$ 的调和函数. 那么函数
\begin{equation}\label{eq:subharmonic-harmonic-lifting}
U(\bm{x})=\begin{cases}\bar u(\bm{x}),&x\in B\\u(\bm{x}),&x\in\Omega\setminus B\end{cases}
\end{equation}
在 $\Omega$ 内也是次调和的 (参见 Problem~\ref{ex:subharmonic-harmonic-lifting}). 称 $U$ 为 $u$ 关于 $B$ 的调和提升.\item 若 $u_1,u_2,\ldots,u_N$ 次调和, 则 $\max\{u_1,u_2,\ldots,u_N\}$ 也次调和.\end{enumerate}
